\documentclass[runningheads]{llncs}

% ---------------------------------------------------------------
% Standalone supplementary material for the main submission.
% The ECCV template makes no provision for an in-document appendix, so
% this material is compiled as a SEPARATE document rather than appended
% to the 14-page main paper.

% TODO REVIEW: keep the same submission ID as the main paper.
% TODO FINAL: comment out the review line, uncomment the plain one.
\usepackage[review,year=2026,ID=*****]{eccv}
%\usepackage{eccv}

\usepackage{eccvabbrv}
\usepackage{graphicx}
\usepackage{booktabs}
\usepackage{bm}
\usepackage{array}
\usepackage{listings}
\usepackage[accsupp]{axessibility}

\graphicspath{{../../thesis/images/}}

\usepackage{hyperref}
\usepackage{orcidlink}


\begin{document}

\title{Supplementary Material\\
\large True Absence or Detection Failure: Generative Augmentation and
Expert-Calibrated Filtering for Long-Tailed Bumblebee Classification}

\titlerunning{Supplementary Material}

\author{Anonymous ECCV submission}
\authorrunning{Anonymous ECCV submission}
\institute{Paper ID \#*****}

\maketitle

\noindent This supplementary material accompanies the main submission.
Section, figure, and table references of the form ``main paper'' point
to the 14-page primary document.


\section{Full Structured Prompt Template}
\label{sec:suppl-prompt}

The complete morphology-guided prompt is assembled programmatically for
each generation call. Placeholders in braces are filled from the
species configuration file. The template comprises six sections: task
framing, subject identity, morphological notes, framing and
composition, pose and viewpoint, and background and quality.

\begin{lstlisting}[basicstyle=\ttfamily\footnotesize,
                   breaklines=true,frame=single,
                   xleftmargin=0.5em,xrightmargin=0.5em]
TASK
Edit the supplied reference photographs into one
new, photorealistic image of {scientific_name}
({common_name}), {caste}, viewed in {view_angle},
on {environment}.

SUBJECT IDENTITY
Genus: Bombus. Subgenus: {subgenus}.
Sex / caste: {caste}.
Body length: {body_length_mm}.
Visual identity must follow the morphological
notes below exactly; do not blend in features
from any other Bombus species.

MORPHOLOGICAL NOTES (front-to-back colour map)
{tergite_colour_map}
{caste_specific_notes}

NEGATIVE CONSTRAINTS (do not produce)
- Bright lemon yellow when the species reference
  says "warm yellow / pale yellow"
- Yellow on the thorax when the species reference
  says BLACK
- Smooth banding when the species reference says
  "abrupt cutoff at T3"
- Visible corbicula on Psithyrus species
  (B. ashtoni, B. flavidus)
- Extra or missing legs, wings, or antennae
- Repetitive scale-like or duplicated body parts
- Visible compression or super-resolution
  artefacts

FRAMING AND COMPOSITION
- Single specimen centered in frame
- Bee occupies 35-55 % of the image area
  (do not zoom in too close)
- {wing_style}
- Sharp focus on diagnostic regions
  (thorax, T1-T6)

POSE AND VIEWPOINT
{view_angle_specific_pose_description}

BACKGROUND AND QUALITY
- Natural {environment} background, shallow
  depth of field, soft daylight
- Photorealistic insect macrophotography style
- No text, no watermark, no border
\end{lstlisting}

The five view angles and their paired wing positions are listed in
\cref{tab:suppl-viewpoints}. Habitat environments are sampled randomly
per generation call rather than cycled deterministically, preventing
spurious environment--species correlations.

\begin{table}[t]
  \caption{View angles and paired wing positions used in the structured
  prompt template.}
  \label{tab:suppl-viewpoints}
  \centering
  \begin{tabular}{@{}ll@{}}
    \toprule
    \textbf{View angle}      & \textbf{Wing style} \\
    \midrule
    Lateral                  & Folded over abdomen at rest \\
    Dorsal                   & Folded over abdomen at rest \\
    Three-quarter anterior   & Slightly spread \\
    Three-quarter posterior  & Folded over abdomen at rest \\
    Frontal                  & Slightly spread \\
    \bottomrule
  \end{tabular}
\end{table}

\section{Volume Ablation}
\label{sec:suppl-volume}

\Cref{fig:suppl-volume} shows ResNet-50 macro-F1 across volumes of
$+50$, $+100$, $+200$, and $+300$ generated images per rare species
under D3 (unfiltered), D4 (LLM-filtered), D5 (centroid-filtered), and
D6 (expert-probe-filtered). All curves remain within a narrow band
overlapping the D1 baseline confidence interval, and rare-species F1
oscillates without a monotonic trend. Adding more generated images of
the same quality does not reliably improve classification; the binding
constraint is selection rather than scale.

\begin{figure}[t]
  \centering
  \includegraphics[width=\linewidth]{ablation/volume_ablation_trends.png}
  \caption{Volume ablation across D3--D6 at $+50$, $+100$, $+200$, and
  $+300$ generated images per rare species, fixed-split evaluation. No
  filter shows a monotonic improvement with volume.}
  \label{fig:suppl-volume}
\end{figure}

\section{Cross-Species Additive Ablation}
\label{sec:suppl-additive}

The subtractive ablation in the main paper removes one rare species'
synthetics at a time. The additive ablation tests the reverse: starting
from the D1 real-only baseline, we add $+200$ D3-unfiltered synthetics
for exactly one rare species (\cref{tab:suppl-additive}). The result is
not purely own-species: adding only \emph{B.~sandersoni} synthetics
yields the highest single-split macro-F1 in the ablation, driven mainly
by a large improvement on \emph{B.~ashtoni} ($+0.218$), while
\emph{B.~sandersoni} itself decreases. Adding only \emph{B.~flavidus}
synthetics helps \emph{B.~flavidus} but harms both other rare species.
Synthetic images therefore reshape decision boundaries across visually
confusable rare classes, not only the class they are labelled as. This
evidence is single-seeded and should be read as directional.

\begin{table}[t]
  \caption{Additive single-species ablation. Each variant adds $+200$
  unfiltered synthetics for exactly one rare species to the D1
  real-only baseline. $\Delta$ values are signed differences from D1.}
  \label{tab:suppl-additive}
  \centering
  \resizebox{\linewidth}{!}{%
  \setlength{\tabcolsep}{3pt}
  \begin{tabular}{@{}lrrrrr@{}}
    \toprule
    Variant
      & Macro F1
      & $\Delta$
      & \emph{ash} F1 ($\Delta$)
      & \emph{san} F1 ($\Delta$)
      & \emph{fla} F1 ($\Delta$) \\
    \midrule
    \textbf{D1\,+\,\emph{san} only} & \textbf{0.857} & \textbf{+0.014}
      & 0.833\,(\textbf{+0.218}) & 0.625\,($-$0.081) & 0.725\,($-$0.022) \\
    D1\,+\,\emph{ash} only          & 0.823 & $-$0.020
      & 0.615\,($+0.000$) & 0.522\,($-$0.184) & 0.667\,($-$0.080) \\
    D1\,+\,\emph{fla} only          & 0.823 & $-$0.020
      & 0.545\,($-$0.070) & 0.375\,($-$0.331) & 0.794\,($+0.048$) \\
    \bottomrule
  \end{tabular}}
\end{table}

\section{Additional Confusion Matrices and Per-Fold Breakdowns}
\label{sec:suppl-confusion}

\Cref{fig:suppl-confusion} shows row-normalised confusion matrices for
the D1 baseline and the D6 expert-probe variant on the fixed test split
(seed 43, the multi-seed median run). The D6 matrix preserves the
dominant diagonal of the baseline while reducing the
\emph{B.~sandersoni}$\to$\emph{B.~vagans} off-diagonal cell that drives
the baseline's rare-tier failure. \Cref{fig:suppl-kfold} reports
per-species F1 deltas relative to D1 across the five stratified
cross-validation folds for D2--D6. Fold-level variance is large for the
rare tier, consistent with small per-fold rare test support, but D6
produces non-negative deltas for at least two of the three rare species
in every fold, while D3, D4, and D5 each include folds where two or
more rare species lose F1.

\begin{figure}[t]
  \centering
  \begin{subfigure}{0.49\linewidth}
    \centering\includegraphics[width=\linewidth]{results/baseline_confusion_matrix.png}
    \caption{D1 baseline.}
  \end{subfigure}
  \hfill
  \begin{subfigure}{0.49\linewidth}
    \centering\includegraphics[width=\linewidth]{results/d6_probe_seed43@f1_confusion_matrix.png}
    \caption{D6 expert-probe.}
  \end{subfigure}
  \caption{Row-normalised confusion matrices on the fixed test split
  (cells with row probability $<\!0.05$ blank for legibility). D6
  reduces the dominant
  \emph{B.~sandersoni}$\to$\emph{B.~vagans} off-diagonal that drives
  the baseline's rare-tier failure.}
  \label{fig:suppl-confusion}
\end{figure}

\begin{figure}[t]
  \centering
  \includegraphics[width=\linewidth]{results/species_f1_delta_kfold.png}
  \caption{Per-species F1 delta relative to D1 across five stratified
  CV folds. Rare species carry most of the augmentation signal; D6
  produces non-negative deltas on at least two of three rare species in
  every fold, while D3, D4, and D5 each include folds where two or more
  rare species lose F1.}
  \label{fig:suppl-kfold}
\end{figure}

\section{Full Pairwise Significance Tests}
\label{sec:suppl-tests}

\Cref{tab:suppl-tests} gives all pairwise paired $t$-tests on fold-level
macro-F1 (five-fold CV) summarised in the main paper. Only the two LLM
comparisons (D1 vs.\ D4 and D2 vs.\ D4) reach significance, both
negative; every other aggregate comparison is non-significant, which is
the basis for the main paper's claim that no generated-filter variant
improves macro-F1 over the baseline or the unfiltered pool.

\begin{table}[t]
  \caption{Pairwise paired $t$-tests on fold-level macro-F1 (5-fold
  CV). $\Delta$ is variant minus reference.}
  \label{tab:suppl-tests}
  \centering
  \begin{tabular}{@{}l rrr l@{}}
    \toprule
    Comparison & $\Delta$ & $t$ & $p$ & Note \\
    \midrule
    D1 vs.\ D2 CNP        & $+0.0051$ & $+1.72$ & $0.16$ & n.s. \\
    D1 vs.\ D3            & $-0.0116$ & $-2.17$ & $0.10$ & marg. \\
    \textbf{D1 vs.\ D4}   & $\bm{-0.0109}$ & $\bm{-2.98}$ & $\bm{0.04}$ & D4 worse \\
    D1 vs.\ D5            & $-0.0082$ & $-1.44$ & $0.22$ & n.s. \\
    D1 vs.\ D6            & $-0.0075$ & $-1.33$ & $0.25$ & n.s. \\
    D2 vs.\ D3            & $-0.0168$ & $-2.59$ & $0.06$ & marg. \\
    \textbf{D2 vs.\ D4}   & $\bm{-0.0160}$ & $\bm{-3.29}$ & $\bm{0.03}$ & D4 worse \\
    D3 vs.\ D6            & $+0.0041$ & $+0.49$ & $0.65$ & n.s. \\
    D5 vs.\ D6            & $+0.0006$ & $+0.09$ & $0.93$ & n.s. \\
    \bottomrule
  \end{tabular}
\end{table}

\end{document}
